\documentclass[12pt]{article}
\usepackage[a4paper,margin=1in]{geometry}
\usepackage{amsmath,amssymb,amsthm}
\usepackage{booktabs}
\usepackage[hidelinks]{hyperref}
\usepackage{setspace}
\doublespacing
\newtheorem{lemma}{Lemma}
\newtheorem{proposition}{Proposition}
\begin{document}
\begin{center}
{\Large\bfseries Supplementary Material}\\
\vspace{4pt}
{\large Directional Evidence on Workforce Reallocation in China's Manufacturing Digital Transformation}
\end{center}

\section*{S1. Proofs of Lemmas 1--3 and Propositions 1--4}
\begin{proof}
The frontier is $\theta(A)=\min\bigl\{\alpha_s,\;\kappa_s^{-1}\bigl(r(A)/w\bigr)\bigr\}$. Since $\kappa_s$ is decreasing and convex, $\kappa_s^{-1}$ is decreasing and convex; since $r$ is decreasing and convex in $A$, $r(A)/w$ is decreasing and convex in $A$. The composition of a decreasing-convex map with a decreasing-convex map is increasing, and concavity follows from the convexity of both maps, giving $\theta'>0$ and $\theta''<0$ on the interior. The limits follow from $\lim_{A\to\infty}r(A)=0$, so $\kappa_s^{-1}(0)=\infty$ and $\theta(A)\to\alpha_s$; normalising the routine share to unit measure yields $\lim_{A\to\infty}\theta(A)=1$.
\end{proof}
\begin{proof}
The retraining condition $w\,p(I_s)\ge w_r+c(I_s)$ defines a threshold over the worker-type distribution; differentiating the resulting share with respect to $I_s$ and applying the envelope theorem with $p'>0$ and $c'<0$ gives $\rho'>0$. Concavity follows from the diminishing returns of the success probability and the convexity of the cost-reduction schedule, standard under a log-concave type distribution; $\rho(0)=0$ because $c(0)$ is prohibitive, and $\rho'(0)>0$ because a small subsidy moves a positive measure of marginal types into retraining.
\end{proof}
\begin{proof}
The productivity gain from automating measure $\theta(A_s)\alpha_s$ of tasks raises the sector's effective TFP. By the CES demand system, derived demand for the surviving non-routine tasks rises with TFP at elasticity $\varepsilon>1$; writing the labour-demand elasticity per unit of automated task share as $\nu_s$ yields~the scale-effect equation of the main text. Then $\partial\sigma_s/\partial A_s=\nu_s\alpha_s L_s\,\theta'(A_s)>0$ by Lemma~Lemma~1 of the main text.
\end{proof}
\begin{proof}
By~the net-change equation of the main text, the sign of $\Delta L_s$ equals the sign of $\nu_s+\rho(I_s)-1$, which is independent of $A_s$ under A1--A3; the sector with capacity above the threshold expands and the one below contracts. When $\nu_M+\rho(I_M)>1>\nu_C+\rho(I_C)$, the two terms of the difference have opposite signs and $\Delta L_M-\Delta L_C>0$, with $\Delta L_M>0$ and $\Delta L_C<0$.
\end{proof}
\begin{proof}
From the decomposition $\Delta N_M=s_{M,0}\Delta N+N_1\Delta s_M$, dividing by $\Delta N_M$ gives $\gamma_M=1-s_{M,0}\Delta N/\Delta N_M$; substituting $\Delta N_M=\Delta L_M$ from~the net-change equation of the main text yields~the composition-ratio equation of the main text, and $\gamma_M>\tfrac12$ rearranges to $\Delta L_M>2\,s_{M,0}\Delta N$.
\end{proof}
\begin{proof}
Differentiating~the cost equation of the main text, A1 gives $\theta'>0$ and A2 gives $\rho'>0$, so $\partial C/\partial A>0$ and the cross-partial is negative.
\end{proof}
\begin{proof}
From~the composition-ratio equation of the main text, $\partial\gamma_M/\partial\Delta L_M=s_{M,0}\Delta N/\Delta L_M^2>0$. By~the net-change equation of the main text, $\partial\Delta L_M/\partial I_M=\theta(A_M)\alpha_M L_M\,\rho'(I_M)>0$ (A2) and $\partial\Delta L_M/\partial\nu_M=\theta(A_M)\alpha_M L_M>0$ (A3); the chain rule gives the two inequalities. Finally $\partial\Delta L_M/\partial A_M=\theta'(A_M)\alpha_M L_M[\nu_M+\rho(I_M)-1]$ has the sign of $\nu_M+\rho(I_M)-1$ (A1).
\end{proof}
\section*{S2. Full Identification Design: Endogeneity and the 101-City Design}

The design specifies the identification strategy for when city-level migrant-worker composition becomes available. A critical identification concern remains: observed manufacturing inflows could reflect either \emph{expansion-driven} reallocation, where growing manufacturing absorbs workers endogenously, or \emph{contraction-driven} reallocation, where construction contraction pushes workers elsewhere regardless of their own dynamics. The two mechanisms imply opposite policy prescriptions. As a first cut we compare two contrasting regimes: Zhejiang, where manufacturing value added expanded at above 8 per cent annual average while construction remained broadly stable, and the three Northeastern provinces together with Shanxi, where manufacturing contracted at approximately $-2$ per cent and construction at approximately $-8$ per cent per year. An event-study model around the August 2023 onset of the first pilot batch is specified with event-time dummies relative to 2023 and province controls for income, urbanisation, industrial structure, higher-education enrolment and robot density. At province level the National Bureau of Statistics does not publish migrant-worker composition, so falsifiable criteria are stated rather than coefficients: under expansion-driven reallocation the Zhejiang-group coefficients are positive and rising while the Northeast group is close to zero, whereas under contraction-driven reallocation the two groups differ little.

The MIIT three-batch pilot (30 cities announced 30 August 2023, 36 on 21 June 2024 and 35 on 30 May 2025) offers a sharper identification opportunity: within a single policy wave, the staggered rollout provides multiple shock timings against which a common-trend assumption can be tested. The identification model is
\[
Y_{c,t} = \alpha + \sum_{b=1}^{3} \beta_b TREAT_{c,b} \times POST_{b,t} + \gamma X_{c,t} + \mu_c + \lambda_t + \varepsilon_{c,t},
\]
where $Y_{c,t}$ is the city-level outcome of interest, $TREAT_{c,b}=1$ if city $c$ belongs to batch $b$ and $POST_{b,t}=1$ at or after the batch rollout year, with standard errors clustered by city. Three extensions are specified ahead of estimation: the primary estimand is the Callaway and Sant'Anna doubly-robust group-time average treatment effect, aggregated with group-size weights; heterogeneity-robust event-study coefficients follow Abraham and Sun (2021) so early-treated cities are not used as controls for late-treated ones; and the pre-trend assumption is tested jointly as $H_0:\beta_{-4}=\beta_{-3}=\beta_{-2}=\beta_{-1}=0$, complemented by a placebo that re-assigns batch weights to fictitious rollout dates drawn uniformly from 2017--2019. Because three outcomes are tested simultaneously, the false discovery rate is controlled with the Benjamini--Hochberg (1995) procedure at $q=0.10$. Two diagnostics address batch assignment: batch membership is regressed on pre-period income, industrial structure, urbanisation and enrolment to expose endogeneity, and the anticipation term $POST_{-1} \times TREAT_{c,b}$ is allowed for because the subsidy caps and the 80 per cent use rule were public before rollout.

The design is declared complete when at least 70 per cent of the 101 cities have outcome data for 2019--2025, the pre-trend test is reported for each batch separately, the three estimands are reported side by side, and the batch-endogeneity and anticipation diagnostics are reported regardless of sign. The 101-city list is used as a coverage and fiscal record; the design specifies the conditions under which it could serve as a basis for control-group construction: the National Bureau of Statistics releases survey reports at national level only, city-level robot density is only partially reported and must be interpolated, and controls such as per-capita GDP and education expenditure must be assembled from city yearbooks. The staggered design relaxes the parallel-trend failure documented in the provincial panel by exploiting multiple treatment timings within a single policy wave, but validity still depends on local pre-trend comparability and batch 1 coincides with the real-estate cycle, so the pre-trend placebo and anticipation test are part of primary reporting rather than optional robustness. Aggregate prefecture-level employment is available, so the binding gap for a causal reading of this paper's outcome is specifically the migrant-worker industry composition, not city-level employment per se.

\section*{S3. Data and reproduction note}

The reproduction files for the international panel (the institutional index components and the automation series) are provided for transparency but should be treated as illustrative. The IFR industrial-robot series is from the licensed (VDMA) database and is not freely reproducible. The local reproduction files were generated with placeholder series where a public source was not available and should not be used to re-estimate the index; the headline IACI values and the panel estimates of the main text are reported as given, not as a re-derivation from these files. Code and data for the Chinese domestic tables (Sections 3--4) are reproducible from the cited public sources.

\section*{References}
\begin{list}{}{\leftmargin=1.5em}
\item Abraham, S. and L. Sun, 2021, Estimating dynamic treatment effects in event studies with heterogeneous treatment effects. \emph{Journal of Econometrics}, 225, pp. 175--99.
\item Benjamini, Y. and Y. Hochberg, 1995, Controlling the false discovery rate: A practical and powerful approach to multiple testing. \emph{Journal of the Royal Statistical Society: Series B}, 57, pp. 289--300.
\item Callaway, B. and P. H. Sant'Anna, 2021, Difference-in-differences with multiple time periods. \emph{Journal of Econometrics}, 225, pp. 200--30.
\end{list}

\end{document}
